\chapter{जोड़: पहले कदम}\label{ch:g1:addition}

तीन लाल कंचे और दो नीले कंचे: सब मिलाकर कितने? मिलाकर पूरी गिनती
करना ही \emph{जोड़ना} है --- पहली संक्रिया, जिसके दो मशहूर चिह्न
हैं $+$ और $=$।

\section{मिलाकर रखना}

\begin{definition}[जोड़]\label{def:g1:addition:def}
\emph{जोड़ना} यानी सब कुछ एक साथ रखकर पूरी गिनती करना। हम
लिखते हैं
\[
3 + 2 = 5
\]
और पढ़ते हैं: ``तीन जमा दो बराबर पाँच''। जोड़ के उत्तर को
\emph{योगफल}\index{योगफल} कहते हैं।
\end{definition}

\begin{omfigure}
\begin{tikzpicture}[scale=0.6]
  \foreach \i in {0,1,2} \fill[omThm] (\i*1.0,0) circle (0.3);
  \node[font=\large] at (3.0,0) {$+$};
  \foreach \i in {0,1} \fill[omDef] (3.8+\i*1.0,0) circle (0.3);
  \node[font=\large] at (5.8,0) {$=$};
  \foreach \i in {0,1,2} \fill[omThm] (6.6+\i*1.0,0) circle (0.3);
  \foreach \i in {0,1} \fill[omDef] (9.6+\i*1.0,0) circle (0.3);
  \node[below, font=\small] at (5.6,-0.7)
    {$3 + 2 = 5$: तीन और दो मिलकर पाँच};
\end{tikzpicture}

{\small चित्रों में जोड़: दोनों समूह मिलकर पाँच का एक ही समूह
बनाते हैं।}
\end{omfigure}

\begin{example}[क्रम से फ़र्क़ नहीं पड़ता]\label{ex:g1:addition:order}
$2 + 3$ और $3 + 2$ वही कंचे मिलाते हैं: दोनों $5$ बनते हैं।
$2 + 9$ निकालने के लिए बड़ी संख्या से शुरू करना तेज़ है: कहो
``नौ'', फिर दो आगे गिनो --- ``दस, ग्यारह''। तो
$2 + 9 = 11$।
\end{example}

\section{दस के साथी}

\begin{example}[10 के पूरक]\label{ex:g1:addition:friends}
जो दो संख्याएँ मिलकर $10$ बनाती हैं, वे ``दस के साथी'' हैं:
\[
1 + 9, \quad 2 + 8, \quad 3 + 7, \quad 4 + 6, \quad 5 + 5 .
\]
दस उँगलियाँ ये सब दिखा देती हैं: $3$ उँगलियाँ उठाओ, तो मुड़ी हुई
$7$ उसकी साथी हैं। दस के साथी याद हों तो कई जोड़ फटाफट हो
जाते हैं।
\end{example}

\begin{omfigure}
\begin{tikzpicture}[scale=0.72]
  % दस-खाना: 5 खानों की 2 पंक्तियाँ
  \draw[thick] (0,0) grid[step=1] (5,2);
  \draw[very thick] (0,0) rectangle (5,2);
  \foreach \i in {0,1,2} \fill[omDef] (\i+0.5,1.5) circle (0.3);
  \foreach \i in {3,4} \draw[omThm, thick] (\i+0.5,1.5) circle (0.3);
  \foreach \i in {0,...,4} \draw[omThm, thick] (\i+0.5,0.5) circle (0.3);
  \node[right, font=\small, align=left] at (5.6,1.0)
    {एक दस-खाना: $10$ खाने।\\
     $3$ भरे, $7$ ख़ाली: $3 + 7 = 10$};
\end{tikzpicture}

{\small दस-खाना दस के हर साथी-जोड़े को एक नज़र में दिखा देता है:
जो भरा है और जो ख़ाली है, वे मिलकर हमेशा $10$ बनते हैं।}
\end{omfigure}

\begin{method}[दस को पार करना]\label{met:g1:addition:crossten}
$8 + 5$ निकालने के लिए:
\begin{enumerate}
  \item $5$ में से उतना लो जितना $8$ को $10$ तक पहुँचने के लिए
        चाहिए: यानी $2$ (उसका दस-साथी);
  \item $8 + 2 = 10$, और $5$ में से $3$ बच रहता है;
  \item $10 + 3 = 13$। तो $8 + 5 = 13$।
\end{enumerate}
$10$ से होकर जाने पर एक कठिन \omterm{def:g1:addition:def}{योगफल} दो आसान योगफलों में बदल जाता है।
\end{method}

\begin{omfigure}
\begin{tikzpicture}[scale=0.62]
  % दस-खाना: 8 नीले, 2 लाल (5 में से लिए हुए)
  \draw[thick] (0,0) grid[step=1] (5,2);
  \draw[very thick] (0,0) rectangle (5,2);
  \foreach \i in {0,...,4} \fill[omDef] (\i+0.5,1.5) circle (0.3);
  \foreach \i in {0,1,2} \fill[omDef] (\i+0.5,0.5) circle (0.3);
  \foreach \i in {3,4} \fill[omThm] (\i+0.5,0.5) circle (0.3);
  % 3 लाल खाने के बाहर बचे
  \foreach \i in {0,1,2} \fill[omThm] (6.2+\i,1.0) circle (0.3);
  \node[below, font=\small] at (2.5,-0.25) {खाना भर गया: $10$};
  \node[below, font=\small] at (7.2,0.55) {$3$ बाहर बचे};
\end{tikzpicture}

{\small दस-खाने पर $8 + 5$: $8$ नीली गोटियाँ $5$ लाल में से $2$
लेकर खाना भर देती हैं --- एक पूरी दहाई और $3$ और, यानी
$13$।}
\end{omfigure}

\begin{omfigure}
\begin{tikzpicture}[scale=0.6]
  \draw[-Stealth] (-0.4,0) -- (15.4,0);
  \foreach \i in {0,...,15} \draw (\i,-0.1) -- (\i,0.1);
  \foreach \i in {0,5,8,10,13,15}
    \node[below, font=\footnotesize] at (\i,-0.12) {\i};
  \draw[-Stealth, omDef, thick] (8,0.3) .. controls (9,0.95) ..
    (10,0.3);
  \node[omDef, font=\small] at (9,1.05) {$+2$};
  \draw[-Stealth, omThm, thick] (10,0.3) .. controls (11.5,1.15) ..
    (13,0.3);
  \node[omThm, font=\small] at (11.5,1.2) {$+3$};
\end{tikzpicture}

{\small संख्या रेखा पर $8 + 5$: पहले $10$ तक छलाँग, फिर बाक़ी।}
\end{omfigure}

\begin{example}[दुगुने]\label{ex:g1:addition:doubles}
दुगुने याद कर लेना अच्छा रहता है:
\[
1+1=2, \quad 2+2=4, \quad 3+3=6, \quad 4+4=8, \quad 5+5=10,
\]
\[
6+6=12, \quad 7+7=14, \quad 8+8=16, \quad 9+9=18, \quad 10+10=20 .
\]
दुगुने के पड़ोसी मुफ़्त में मिलते हैं: $6 + 7$, $6 + 6$ से एक ज़्यादा है, यानी $13$।
\end{example}

\section{छोटी-छोटी समस्याएँ}

\begin{example}\label{ex:g1:addition:problem}
``लेआ के पास $7$ स्टिकर हैं। उसे $6$ और मिलते हैं। अब कितने?''
और मिलना यानी जोड़ना: $7 + 6 = 13$ (दस से होकर:
$7 + 3 = 10$, फिर $+3$)। लेआ के पास $13$ स्टिकर हैं। जवाब हमेशा
एक छोटे वाक्य में पूरा करो!
\end{example}

\section{अभ्यास}

\begin{exercise}[$\star$]\label{exo:g1:addition:1}
कंचे बनाकर जोड़ो: $4 + 3$; \quad $5 + 2$; \quad
$6 + 4$।
\end{exercise}

\begin{exercise}[$\star$]\label{exo:g1:addition:2}
बड़ी संख्या से शुरू करके जोड़ो: $3 + 8$; \quad
$2 + 12$; \quad $4 + 9$।
\end{exercise}

\begin{exercise}[$\star$]\label{exo:g1:addition:3}
इनके दस-साथी बताओ: $7$; \ $4$; \ $9$; \ $5$; \ $1$।
\end{exercise}

\begin{exercise}[$\star$]\label{exo:g1:addition:4}
दस से होकर जोड़ो (\cref{met:g1:addition:crossten}):
$9 + 4$; \quad $8 + 7$; \quad $6 + 5$।
\end{exercise}

\begin{exercise}[$\star$]\label{exo:g1:addition:5}
दुगुने बोलो, फिर उनकी मदद से जोड़ो: $7 + 8$; \quad
$5 + 6$; \quad $9 + 8$।
\end{exercise}

\begin{exercise}[$\star$]\label{exo:g1:addition:6}
ख़ाली जगह भरो:
\[
6 + \;?\; = 10, \qquad
\;?\; + 5 = 12, \qquad
10 + \;?\; = 17 .
\]
\end{exercise}

\begin{exercise}[$\star$]\label{exo:g1:addition:7}
``बस में $8$ बच्चे हैं; $5$ बड़े चढ़ते हैं। बस में कुल कितने लोग
हैं?'' जोड़ और जवाब का वाक्य लिखो।
\end{exercise}

\begin{exercise}[$\star$]\label{exo:g1:addition:8}
टॉम के पास $6$ लाल कारें और $7$ नीली कारें हैं। मिया के पास $10$
कारें हैं। किसके पास ज़्यादा कारें हैं?
\end{exercise}

\begin{exercise}[$\star\star$]\label{exo:g1:addition:9}
तीन पासों पर $4$, $6$ और $5$ आए हैं। कुल कितना? (पहले दो पासे
जोड़ो --- वे दो चुनो जो दस के साथी बनते हैं!)
\end{exercise}
